\section{Results status}
\label{sec:results}

ORION Paper III has deliberately separate evidence layers. The original
end-to-end protocol \texttt{P3.cross-domain-atlas.v1} remains \conststatus{} for
real-text extraction, broad construct validity, recoverability and downstream
utility. A narrower public-reference mapping protocol has been executed using
independently produced human/expert resources, followed by a disjoint
execution-frozen confirmatory replication. The numbers below come only from
immutable archived artifacts.

\subsection{Frozen public-reference mapping atlas}

The public-reference builder drew only from pinned external authorities: MUSE
expert annotations, SciFact expert SUPPORT/CONTRADICT annotations, and versioned
SciSchema expert schemas. Unsupported semantic coordinates were not filled with
model guesses. The first deterministic build produced 32 cases across five
recorded strata and three case families: different-name/same-referent,
polarity/modality/attribution/context, and valid representation mappings. The
portable gold artifact has SHA-256
\texttt{35f9e39b75ff53b7f0ec82cd03ebcaaa82509ee0aea3f5b96aac3fd62c854ed8}.
A second freeze pass reproduced it byte-for-byte.

This atlas evaluates the already-structured projection-to-mapping calculus. It
is not an evaluation of a language model extracting projections from raw
papers.

\subsection{Initial public-reference mapping result}

On the first 32 frozen mapping cases, the ORION semantic comparison rule
classified all 32 expected relations correctly (accuracy 1.000), with
false-merge rate 0.000 and false-split rate 0.000. A flat
predicate-canonicalization control had accuracy 0.875 and false-merge rate
0.125, with no false splits. The paired ORION-minus-flat false-merge difference
was $-0.125$ with a 95\% paired percentile-bootstrap interval of
$[-0.250,-0.03125]$ (10,000 resamples; bootstrap seed 20260817).

The conservative exact-coordinate control also produced zero false merges but
abstained on 0.125 of the cases and achieved accuracy 0.875. The initial result
therefore did not reduce to ``more conservative is better'': the typed rule
retained the accepted mappings while the exact-match control left four cases
unresolved.

\subsection{Prospectively frozen confirmatory replication}

Because the first public-reference run was design-frozen but its execution
bindings were not yet fully bound, we treated it as an initial result rather
than the final confirmatory authority. We then created a second, disjoint
32-case holdout using the same outcome-blind round-robin construction at
selection offset 32. The holdout was frozen before any confirmatory ORION or
control output existed, with zero case overlap with the initial atlas. Its
portable SHA-256 is
\texttt{13a76c68c149c2552f3543babeca6e1ad5afe23c45ea9c0dc365c1445cf2782b}.
The execution manifest \texttt{P3.public-reference-mapping.v1.1-confirmatory}
bound that hash, the upstream source revisions, evaluator Git blobs, bootstrap
seed, practical margins, and pass rule before evaluation.

The predeclared confirmatory rule required both (i) ORION-minus-flat
false-merge difference at most $-0.05$ with the upper end of its paired 95\%
bootstrap interval below zero, and (ii) ORION-minus-exact false-split difference
no larger than $+0.03$ with the interval upper bound no larger than $+0.03$.
The disjoint holdout passed both conditions. ORION again had false-merge rate
0.000 and false-split rate 0.000. Flat predicate canonicalization had
false-merge rate 0.1875 and accuracy 0.8125. The paired ORION-minus-flat
false-merge difference was $-0.1875$ with 95\% interval
$[-0.34375,-0.0625]$. The ORION-minus-exact false-split difference was 0.000
with interval $[0.000,0.000]$. The confirmatory primary verdict was therefore
\textbf{PASS} under the rule frozen before those outputs were generated.

The effect is concentrated where the holdout actually has discriminating
support. Among 13 polarity/modality/attribution/context cases, flat
canonicalization false-merged 6/13 (0.4615), while ORION false-merged none. The
13 different-name/same-referent cases and six representation-mapping cases were
correct under all compared rules. These strata are reported descriptively and
are not pooled with the original 32 cases to manufacture a larger confirmatory
sample.

\subsection{Covered ablations}

On the confirmatory holdout, forcing compatibility without the
obstruction/contradiction distinction increased false merges by $+0.1875$
(95\% paired bootstrap interval $[+0.0625,+0.34375]$) and reduced accuracy by
the same absolute amount. Removing the modality/polarity/attribution/discourse
coordinate family produced the same false-merge increase on the covered
SciFact controls. Other implemented coordinate-removal ablations have zero
measured effect in this atlas; that is a coverage result, not evidence that
those coordinates are unnecessary. The confirmatory protocol treats these
ablations as descriptive secondary analyses rather than significance claims.

\subsection{Verified local exact-world semantics}

Separate synthetic exact-world tests continue to verify implementation
semantics for referent and construct identity, measurement equivalence, context,
\glue{}, \obstruction{}, \pluralview{}, preservation conditions and typed
mapping hypotheses. These tests are engineering/semantic falsifiers and carry no
external empirical authority.

\subsection{What remains \texttt{CANNOT\_CHECK}}

The replicated public-reference route still does not establish the original
paper-wide primary claim. The following remain unverified under
\texttt{P3.cross-domain-atlas.v1}:

\begin{itemize}
  \item \textbf{End-to-end H1:} whether a system starting from raw scientific
        text reduces false merges and false contradictions against the strongest
        resource-matched synthesis/schema baselines across the full frozen case
        family set.
  \item \textbf{H2 recoverability:} whether readers can recover source
        projection, mapping assumptions and non-preserved information from
        generated global statements.
  \item \textbf{H3 downstream obstruction value:} whether retaining incompatible
        views improves scientific answering or decisions relative to forced
        canonicalization.
  \item \textbf{Full H4 coordinate necessity:} the public-reference atlases
        directly stress only a subset of the frozen coordinate families; zero
        ablation effects outside that support are not evidence of dispensability.
\end{itemize}

A text-to-scientific-meaning provider port and structured projection semantics
exist, but no prospectively bound raw-text extractor/provider has been executed
for the headline end-to-end protocol. Future results must be added only from
newly frozen artifacts; neither public-reference run may be generalized beyond
the mapping scope actually evaluated.

\endinput
